% problems9.tex — Problem Set for Chapter 9

\begin{problemset}

\subsection*{Analytical Exercises}

\exercise[1]
Prove Proposition~9.2(b). \textbf{Hint:} By the definition of $\Delta\xi_t$, we have $\xi_t = \xi_{t-1} +
\Delta\xi_t$. So
\[
(\xi_t)^2 = (\xi_{t-1} + \Delta\xi_t)^2.
\]
From this, derive
\[
\Delta\xi_t \cdot \xi_{t-1} = \bigl(\tfrac{1}{2}\bigr)\bigl[(\xi_t)^2 - (\xi_{t-1})^2 - (\Delta\xi_t)^2\bigr].
\]
Sum this over $t = 1, 2, \ldots, T$ to obtain
\[
\sum_{t=1}^{T}\Delta\xi_t \cdot \xi_{t-1}
= \bigl(\tfrac{1}{2}\bigr)\bigl[(\xi_T)^2 - (\xi_0)^2\bigr]
- \bigl(\tfrac{1}{2}\bigr)\sum_{t=1}^{T}(\Delta\xi_t)^2.
\]
Divide both sides by $T$ to obtain
\[
\frac{1}{T}\sum_{t=1}^{T}\Delta\xi_t \cdot \xi_{t-1}
= \frac{1}{2}\left(\frac{\xi_T}{\sqrt{T}}\right)^2
- \frac{1}{2}\left(\frac{\xi_0}{\sqrt{T}}\right)^2
- \frac{1}{2T}\sum_{t=1}^{T}(\Delta\xi_t)^2. \quad (*)
\]
Apply Proposition~6.9 and the Ergodic Theorem to the terms on the right-hand
side of $(*)$.

\exercise[2]
(Proof of Proposition~9.4) \enspace Let $\hat{\rho}^\mu$ be the OLS estimate of $\rho$ in the AR(1)
equation with intercept,~(9.3.11), and define the demeaned series $\{y_{t-1}^\mu\}$ $(t =
1, 2, \ldots, T)$ by
\[
y_{t-1}^\mu \equiv y_{t-1} - \bar{y}, \quad
\bar{y} \equiv \frac{y_0 + y_1 + \cdots + y_{T-1}}{T}
\quad (t = 1, 2, \ldots, T). \qquad (1)
\]
This $y_{t-1}^\mu$ equals the OLS residual from the regression of $y_{t-1}$ on a constant for
$t = 1, 2, \ldots, T$.

\begin{enumerate}[label=(\alph*)]
\item Derive
\begin{equation*}
\hat{\rho}^\mu - 1 = \frac{\sum_{t=1}^{T}\Delta y_t\, y_{t-1}^\mu}{\sum_{t=1}^{T}(y_{t-1}^\mu)^2}.
\tag{2}
\end{equation*}
\textbf{Hint:} By the Frisch-Waugh Theorem, $\hat{\rho}^\mu$ is numerically equal to the OLS
coefficient estimate in the regression of $y_t$ on $y_{t-1}^\mu$ (\emph{without} intercept). Thus,
\[
\hat{\rho}^\mu = \frac{\sum_{t=1}^{T} y_t\, y_{t-1}^\mu}{\sum_{t=1}^{T}(y_{t-1}^\mu)^2}
\quad \text{or} \quad
\hat{\rho}^\mu - 1 = \frac{\sum_{t=1}^{T}(y_t - y_{t-1}) \cdot y_{t-1}^\mu}{\sum_{t=1}^{T}(y_{t-1}^\mu)^2}.
\tag{3}
\]
Use the fact that $\sum_{t=1}^{T} y_{t-1}^\mu = 0$ to show that $\sum_{t=1}^{T}(y_t - y_{t-1}) \cdot y_{t-1}^\mu =
\sum_{t=1}^{T}(y_t - y_{t-1}) \cdot y_{t-1}^\mu$.

\item (Easy) \enspace Show that, under the null hypothesis that $\{y_t\}$ is driftless I(1) (not
necessarily driftless random walk),
\[
T \cdot (\hat{\rho}^\mu - 1)
\;\dto\;
\frac{\frac{\lambda^2}{2}\big([W^\mu(1)]^2 - [W^\mu(0)]^2\big) - \frac{\gamma_0}{2}}{\lambda^2 \cdot \int_0^1 [W^\mu(r)]^2\,\mathrm{d}r},
\tag{4}
\]
where $\gamma_0$ is the variance of $\Delta y_t$, $\lambda^2$ is the long-run variance of $\Delta y_t$, and
$W^\mu(\cdot)$ is a demeaned standard Brownian motion introduced in Section~9.2.
\textbf{Hint:} Use Proposition~9.2(c) and~(d) with $\xi_t = y_t$.

\item (Trivial) \enspace Verify that $T \cdot (\hat{\rho}^\mu - 1)$ converges to $DF_\rho^\mu$ (defined in~(9.3.13)) if
$\{y_t\}$ is a driftless random walk. \textbf{Hint:} If $y_t$ is a random walk, then $\lambda^2 = \gamma_0$.

\item (Optional) \enspace Let $s$ be the standard error of regression for the AR(1) equation
with intercept,~(9.3.11). Show that $s^2$ is consistent for $\gamma_0$ ($\equiv \Var(\Delta y_t)$)
under the null that $\{y_t\}$ is driftless I(1) (so $\Delta y_t$ is zero-mean I(0) satisfying
(9.2.1)--(9.2.3)). \textbf{Hint:} Let $(\hat{\alpha}^*, \hat{\rho}^\mu)$ be the OLS estimates of $(\alpha^*, \rho)$. Show
that $\hat{\alpha}^*$ converges to zero in probability. Write $s^2$ as
\begin{align*}
s^2 &= \frac{1}{T-1}\sum_{t=1}^{T}
\bigl[(\Delta y_t - \hat{\alpha}^*) - (\hat{\rho}^\mu - 1)y_{t-1}\bigr]^2 \\
&= \frac{1}{T-1}\sum_{t=1}^{T}(\Delta y_t - \hat{\alpha}^*)^2 \\
&\quad - \frac{2}{T-1}\cdot [T \cdot (\hat{\rho}^\mu - 1)] \cdot
\frac{1}{T}\sum_{t=1}^{T}(\Delta y_t - \hat{\alpha}^*) \cdot y_{t-1} \\
&\quad + \frac{1}{T-1}\cdot [T \cdot (\hat{\rho}^\mu - 1)]^2 \cdot
\frac{1}{T^2}\sum_{t=1}^{T}(y_{t-1})^2.
\tag{5}
\end{align*}
Since $\plim\,\hat{\alpha}^* = 0$, it should be easy to show that the first term on the right-hand side of~(5) converges in probability to $\gamma_0 = \E[(\Delta y_t)^2]$. To prove that
the second term converges to zero in probability, take for granted that
\[
\frac{1}{\sqrt{T}}\frac{1}{T}\sum_{t=1}^{T}y_{t-1}
\;\dto\; \lambda \int_0^1 W(r)\,\mathrm{d}r.
\tag{6}
\]

\item Show that $t^\mu \dto DF_t^\mu$ if $\{y_t\}$ is a driftless random walk. \textbf{Hint:} Since
$\{y_t\}$ is a driftless random walk, $\lambda^2 = \gamma_0 = \sigma^2$ ($\equiv$ variance of $\varepsilon_t$). So $s$ is
consistent for $\sigma$.
\[
t^\mu = \frac{\hat{\rho}^\mu - 1}{s \times \sqrt{(2,2) \text{ element of } (\mathbf{X}'\mathbf{X})^{-1}}}.
\tag{7}
\]
Show that the square root of the $(2, 2)$ element of $\bigl(\sum_{t=1}^{T}\mathbf{x}_t\mathbf{x}_t'\bigr)^{-1}$ is
\[
\frac{1}{\sqrt{\sum_{t=1}^{T}(y_{t-1}^\mu)^2}},
\]
where $\mathbf{x}_t \equiv [1, y_{t-1}]'$.
\end{enumerate}

\exercise[3]
(Proof of Proposition~9.5) \enspace Let $\hat{\rho}^\tau$ be the OLS estimate of $\rho$ in the AR(1)
equation with time trend, (9.3.17), and define the detrended series $\{y_{t-1}^\tau\}$ $(t =
1, 2, \ldots, T)$ by
\[
y_{t-1}^\tau \equiv y_{t-1} - \hat{\alpha} - \hat{\delta} \cdot t,
\tag{8}
\]
where $(\hat{\alpha}, \hat{\delta})$ are the OLS coefficient estimates for the regression of $y_{t-1}$ on
$(1, t)$ for $t = 1, 2, \ldots, T$.

\begin{enumerate}[label=(\alph*)]
\item Derive
\[
\hat{\rho}^\tau - 1 = \frac{\sum_{t=1}^{T}\Delta y_t\, y_{t-1}^\tau}{\sum_{t=1}^{T}(y_{t-1}^\tau)^2}.
\tag{9}
\]
\textbf{Hint:} By the Frisch-Waugh Theorem, $\hat{\rho}^\tau$ is numerically equal to the OLS
estimate of the $y_{t-1}^\tau$ coefficient in the regression of $y_t$ on $y_{t-1}^\tau$ (\emph{without} intercept or time). Thus,
\[
\hat{\rho}^\tau = \frac{\sum_{t=1}^{T} y_t\, y_{t-1}^\tau}{\sum_{t=1}^{T}(y_{t-1}^\tau)^2}
\quad \text{or} \quad
\hat{\rho}^\tau - 1 = \frac{\sum_{t=1}^{T}(y_t - y_{t-1}) \cdot y_{t-1}^\tau}{\sum_{t=1}^{T}(y_{t-1}^\tau)^2}.
\tag{10}
\]
Show that $\sum_{t=1}^{T}(y_t - y_{t-1}) \cdot y_{t-1}^\tau = \sum_{t=1}^{T}(y_t - y_{t-1}) \cdot y_{t-1}^\tau$. (By construction,
$\sum_{t=1}^{T} y_{t-1}^\tau = 0$ and $\sum_{t=1}^{T} t \cdot y_{t-1}^\tau = 0$ (these are the normal equations).)

\item (Easy) \enspace Show that, under the assumption that $y_t = \alpha + \delta \cdot t + \xi_t$ where $\{\xi_t\}$
is driftless I(1),
\[
T \cdot (\hat{\rho}^\tau - 1) \;\dto\;
\frac{\frac{\lambda^2}{2}\big([W^\tau(1)]^2 - [W^\tau(0)]^2\big) - \frac{\gamma_0}{2}}{\lambda^2 \int_0^1 [W^\tau(r)]^2\,\mathrm{d}r},
\tag{11}
\]
where $\gamma_0$ is the variance of $\Delta y_t$, $\lambda^2$ is the long-run variance of $\Delta y_t$, and
$W^\tau(\cdot)$ is a detrended standard Brownian motion introduced in Section~9.2.
\textbf{Hint:} Let $\xi_{t-1}^\tau$ be the residual from the hypothetical regression of $\xi_{t-1}$ on
$(1, t)$ for $t = 1, 2, \ldots, T$ where $\xi_t \equiv y_t - \alpha - \delta \cdot t$ (this regression is
hypothetical because we do not observe $\{\xi_t\}$). Then $\xi_{t-1}^\tau = y_{t-1}^\tau$ for $t = 1,
2, \ldots, T$, because $\xi_t$ differs from $y_t$ only by a linear trend $\alpha + \delta \cdot t$. Use
Proposition~9.2(e) and~(f). Since $\Delta y_t = \delta + \Delta\xi_t$ under the null, the variance
of $\Delta\xi_t$ equals that of $\Delta y_t$ and the long-run variance of $\Delta\xi_t$ equals that of
$\Delta y_t$.

\item (Trivial) \enspace Verify that $T \cdot (\hat{\rho}^\tau - 1)$ converges to $DF_\rho^\tau$ (defined in~(9.3.19))
if $\{y_t\}$ is a random walk with or without drift.
\end{enumerate}

\exercise[4]
(Proof of~(9.4.16)) \enspace In the text,
\[
\frac{1}{T}\sum_{t=1}^{T} y_{t-1}\,\varepsilon_t
\;\dto\; c_1 \equiv \tfrac{1}{2}\sigma^2 \cdot \psi(1) \cdot [W(1)^2 - 1]
\tag{9.4.16}
\]
was left unproven. Prove this under the condition that $\{y_t\}$ is driftless I(1)
so that $\{\Delta y_t\}$ is zero-mean I(0) satisfying (9.2.1)--(9.2.3). \textbf{Hint:} The $\varepsilon_t$ in
(9.4.16) is the innovation process in the MA representation of $\{\Delta y_t\}$. Using
the Beveridge-Nelson decomposition, we have
\[
y_t = \psi(1)w_t + \eta_t + (y_0 - \eta_0), \quad w_t \equiv \varepsilon_1 + \varepsilon_2 + \cdots + \varepsilon_t.
\]
So
\[
\frac{1}{T}\sum_{t=1}^{T} y_{t-1}\,\varepsilon_t
= \psi(1)\frac{1}{T}\sum_{t=1}^{T} w_{t-1}\,\varepsilon_t
+ \frac{1}{T}\sum_{t=1}^{T}\eta_{t-1}\,\varepsilon_t
+ (y_0 - \eta_0)\frac{1}{T}\sum_{t=1}^{T}\varepsilon_t.
\tag{$*$}
\]
Show that the second and the third term on the right-hand side of this equation converge in probability to zero. Once you have done Analytical Exercise~1
above, it should be easy to show that
\[
\frac{1}{T}\sum_{t=1}^{T} w_{t-1}\,\varepsilon_t
\;\dto\; \left(\frac{\sigma^2}{2}\right)[W(1)^2 - 1].
\]

\exercise[5]
(Optional, proof of Proposition~9.7) \enspace To simplify the calculation, set $p = 1$, so
the augmented autoregression is
\[
y_t = \alpha^* + \rho\, y_{t-1} + \zeta_1\,\Delta y_{t-1} + \varepsilon_t.
\]
Let $(\hat{\alpha}^*, \hat{\rho}^\mu, \hat{\zeta}_1)$ be the OLS coefficient estimates. Show that
\[
\frac{T \cdot (\hat{\rho}^\mu - 1)}{1 - \hat{\zeta}_1}
\;\dto\; DF_\rho^\mu.
\]
\textbf{Hint:} By the Frisch-Waugh Theorem, $(\hat{\rho}^\mu, \hat{\zeta}_1)$ are numerically the same as the
OLS coefficient estimates from the regression of $y_t$ on $(y_{t-1}^\mu, (\Delta y_{t-1})^{(\mu)})$, where
$y_{t-1}^\mu$ is the residual from the regression of $y_{t-1}$ on a constant and $(\Delta y_{t-1})^{(\mu)}$ is
the residual from the regression of $\Delta y_{t-1}$ on a constant for $t = 1, 2, \ldots, T$.
Therefore, if $\mathbf{b} = (\hat{\rho}^\mu, \hat{\zeta}_1)'$, $\bm{\beta} = (\rho, \zeta_1)'$, and $\mathbf{x}_t = (y_{t-1}^\mu, (\Delta y_{t-1})^{(\mu)})'$, then
they satisfy~(9.4.8). So

\begin{align*}
&(1, 1) \text{ element of } \mathbf{A}_T = \frac{1}{T^2}\sum_{t=1}^{T}(y_{t-1}^\mu)^2, \\
&\text{1st element of } \mathbf{c}_T = \frac{1}{T}\sum_{t=1}^{T}y_{t-1}^\mu \cdot \varepsilon_t.
\end{align*}

Apply Proposition~9.2(c) and~(d) with $\xi_t = y_t$ to these expressions.

\exercise[6]
(Phillips tests) \enspace Suppose that $\{y_t\}$ is driftless I(1), so $\Delta y_t$ can be serially correlated. Nevertheless, estimate the AR(1) equation without intercept, rather than
the augmented autoregression without intercept, on $\{y_t\}$. Let $\hat{\rho}$ be the OLS
estimate of the $y_{t-1}$ coefficient. Define
\begin{align*}
Z_\rho &\equiv T \cdot (\hat{\rho} - 1)
- \frac{1}{2} \cdot \left(\frac{T^2 \cdot \hat{\sigma}_{\hat{\rho}}^2}{s^2}\right) \cdot (\hat{\lambda}^2 - \hat{\gamma}_0), \\
Z_t &\equiv \frac{s}{\hat{\lambda}} \cdot t
- \frac{1}{2} \cdot \left(\frac{T \cdot \hat{\sigma}_{\hat{\rho}}}{s \cdot \hat{\lambda}}\right) \cdot (\hat{\lambda}^2 - \hat{\gamma}_0),
\end{align*}
where $\hat{\sigma}_{\hat{\rho}}$ is the standard error of $\hat{\rho}$, $s$ is the OLS standard error of regression,
$\hat{\lambda}^2$ is a consistent estimate of $\lambda^2$, $\hat{\gamma}_0$ is a consistent estimate of $\gamma_0$, and $t$ is the
$t$-value for the hypothesis that $\rho = 1$.

\begin{enumerate}[label=(\alph*)]
\item Show that $Z_\rho$ and $Z_t$ can be written as
\begin{align*}
Z_\rho &= \frac{\frac{1}{T}\sum_{t=1}^{T}\Delta y_t\, y_{t-1}
- \frac{1}{2}(\hat{\lambda}^2 - \hat{\gamma}_0)}{\frac{1}{T^2}\sum_{t=1}^{T}(y_{t-1})^2}, \\
Z_t &= \frac{s}{\hat{\lambda}} \cdot t
- \frac{\frac{1}{2}(\hat{\lambda}^2 - \hat{\gamma}_0)}{\hat{\lambda} \times \sqrt{\frac{1}{T^2}\sum_{t=1}^{T} y_{t-1}^2}}.
\end{align*}
\textbf{Hint:} By the algebra of least squares,
\[
\frac{\hat{\sigma}_{\hat{\rho}}}{s} = \frac{1}{\sqrt{\sum_{t=1}^{T}(y_{t-1})^2}}.
\]

\item Show that
\[
Z_\rho \;\dto\; DF_\rho \quad \text{and} \quad Z_t \;\dto\; DF_t.
\]
\textbf{Hint:} Use Proposition~9.2(a) and~(b) with $\xi_t = y_t$.
\end{enumerate}

\exercise[7]
(Optional) \enspace Show that $\alpha(L)$ in~(9.2.5) is absolutely summable. \textbf{Hint:} A one-line
proof is
\[
\sum_{j=0}^{\infty}|\alpha_j|
= \sum_{j=0}^{\infty}\left|\;-\sum_{i=j+1}^{\infty}\psi_i\;\right|
\leq \sum_{j=0}^{\infty}\sum_{i=j+1}^{\infty}|\psi_i|
= \sum_{i=0}^{\infty} i\,|\psi_i| < \infty.
\]
Justify each of the equalities and inequalities.

\subsection*{Monte Carlo Exercises}

\exercise[1]
(You have less power with time) \enspace In this simulation, we wish to verify that the
finite-sample power against stationary alternatives is greater with the DF $t^\mu$-test
than with the DF $t^\tau$-test, that is, inclusion of time in the AR(1) equation reduces
power. The DGP we examine as the stationary alternative is a stationary AR(1)
process:
\[
y_t = \rho y_{t-1} + \varepsilon_t, \quad \rho = 0.95, \quad
\varepsilon_t \sim \text{i.i.d.\ } N(0, 1), \quad y_0 = 0.
\]
Choose the sample size $T$ to be 100 and the nominal size to be 5 percent.
Your computer program for calculating the finite-sample power should have the
following steps.
\begin{enumerate}
\item Set counters to zero. There should be two counters, one for the DF $t^\mu$ and
the other for $t^\tau$. They will record the number of times the (false) null gets
rejected.

\item Start a do loop of a large number of replications (1 million, say). In each
replication, do the following.

\begin{enumerate}[label=(\roman*)]
\item Generate the data vector $(y_0, y_1, \ldots, y_T)$.

\textit{Gauss Tip:} As was mentioned in the Monte Carlo exercise for Chapter~1, to generate $(y_1, y_2, \ldots, y_T)$, avoid creating a do loop within the
current do loop; use the matrix operators. The matrix formula is
\[
\underset{(T \times 1)}{\mathbf{y}}
= \underset{(T \times 1)}{\mathbf{r}} \cdot y_0
+ \underset{(T \times T)(T \times 1)}{\mathbf{A}\;\bm{\varepsilon}}
\]
where
\[
\mathbf{y} = \begin{bmatrix} y_1 \\ y_2 \\ \vdots \\ y_T \end{bmatrix}, \quad
\mathbf{r} = \begin{bmatrix} \rho \\ \rho^2 \\ \vdots \\ \rho^T \end{bmatrix}, \quad
\mathbf{A} = \begin{bmatrix}
1 & 0 & \cdots & 0 \\
\rho & 1 & \cdots & 0 \\
\rho^2 & \rho & \ddots & \vdots \\
\vdots & \vdots & \ddots & 1 \\
\rho^{T-1} & \rho^{T-2} & \cdots & \rho
\end{bmatrix}, \quad
\bm{\varepsilon} = \begin{bmatrix} \varepsilon_1 \\ \varepsilon_2 \\ \vdots \\ \varepsilon_T \end{bmatrix}.
\]

\item Calculate $t^\mu$ and $t^\tau$ from $(y_0, y_1, \ldots, y_T)$. (So the sample size is actually $T + 1$.)

\item Increase the counter for $t^\mu$ by~1 if $t^\mu < -2.86$ (the 5 percent critical
value for $DF_t^\mu$). Increase the counter for $t^\tau$ by~1 if $t^\tau < -3.41$ (the 5
percent critical value for $DF_t^\tau$).
\end{enumerate}

\item After the do loop, for each statistic, divide the counter by the number of
replications to calculate the frequency of rejecting the null. This frequency
converges to the finite-sample power as the number of Monte Carlo replications goes to infinity.
\end{enumerate}

Power should be about 0.124 for the DF $t^\mu$-test (so only about 12 percent of
the time can we reject the false null!) and 0.092 for the DF $t^\tau$-test. So indeed,
power is less when time is included in the regression, at least against this particular DGP.

\exercise[2]
(Size distortions) \enspace In this Monte Carlo simulation, we calculate the size distortions for the ADF $\rho^\mu$- and $t^\mu$-tests. Our null DGP is the one examined by
Schwert (1989):
\[
y_t = \rho y_{t-1} + \varepsilon_t + \theta\varepsilon_{t-1}, \quad
\rho = 1, \quad \theta = -0.8,
\]
\[
\varepsilon_t \sim \text{i.i.d.\ } N(0, 1), \quad y_0 = 0, \quad \varepsilon_0 = 0.
\]
For $T = 100$ and the nominal size of 5 percent, calculate the finite-sample or
exact size for $\hat{\rho}^\mu$ and $t^\mu$. It would be useful to select the number of lagged
changes in the augmented autoregression by one of the data-dependent methods (the sequential $t$, the AIC, or the BIC), but to keep the computer program
simple and also to save CPU time, set it equal to~4. Thus, the regression equation, as opposed to the DGP, is
\[
y_t = \text{const.} + \rho y_{t-1} + \zeta_1 \cdot (y_{t-1} - y_{t-2})
+ \zeta_2 \cdot (y_{t-2} - y_{t-3})
\]
\[
\quad + \zeta_3 \cdot (y_{t-3} - y_{t-4}) + \zeta_4 \cdot (y_{t-4} - y_{t-5}) + \text{error}.
\]
Since the sample includes $y_0$, the sample period that can accommodate four
lagged changes is from $t = 5$ (not $p + 2 = 6$) to $t = T$ and the actual sample
size is $T - 4$. Verify that the size distortion is less for the ADF $t^\mu$- than for the
ADF $\rho^\mu$-test. (The size should be about 0.496 for the ADF $\rho^\mu$ and 0.290 for
the ADF $t^\mu$.)

\subsection*{Empirical Exercises}

Read Lothian and Taylor (1996, pp.~488--509) (but skip over their discussion of the
Phillips-Perron tests) before answering. The file LT.ASC has annual data on the
following:

\medskip\noindent
Column 1: year\\
Column 2: dollar/sterling exchange rate (call it $S_t$)\\
Column 3: U.S.\ WPI (wholesale price index), normalized to 100 for 1914 ($P_t$)\\
Column 4: U.K.\ WPI, normalized to 100 for 1914 ($P_t^*$).

\medskip\noindent
The sample period is 1791 to 1990 (200 observations). These are the same data
used by Lothian and Taylor (1996) and were made available to us. For data sources
and how the authors combined them to construct consistent time series, see the
Appendix of Lothian and Taylor (1996). (According to the Appendix, the exchange
rate (and probably WPIs) are annual averages. This is rather unfortunate and point-in-time data should have been used, because of the time aggregation problem: if a
variable follows a random walk on a daily basis, the annual averages of the variable
do not follow a random walk. See part~(f) below for more on this.)

Calculate the dollar/sterling real exchange rate as
\[
z_t \equiv \ln(S_t) - \ln(P_t) + \ln(P_t^*).
\tag{1}
\]
Since $S_t$ is in dollars per pound, an increase means sterling appreciation. The plot
of the real exchange rate in Figure~9.4 shows no clear time trend, so we apply the
ADF $t^\mu$-test. Therefore, the augmented autoregression to be estimated is
\[
z_t = \text{const.} + \rho z_{t-1} + \zeta_1 \cdot (z_{t-1} - z_{t-2}) + \cdots
+ \zeta_p \cdot (z_{t-p} - z_{t-p-1}) + \text{error},
\tag{2}
\]
where $p$ is the number of lagged changes included in the augmented autoregression
(which was denoted $\hat{p}$ in the text).

\begin{enumerate}[label=(\alph*)]
\item (ADF with automatic lag selection) \enspace For the entire sample period of 1791--1990, apply the sequential $t$-rule, the AIC, and the BIC to select $p$ (the number of lagged changes to be included). Set $p_{\max}$ (the maximum number of
lagged changes in the augmented autoregression, denoted $p_{\max}(T)$ in the text) by~(9.4.31) (so $p_{\max} = 14$). For the sequential $t$ test, use the significance level
of 10 percent. (You should find $p = 14$ by the sequential $t$, 14 by AIC, and 0
by BIC.) Verify that the value of the $t^\mu$-statistic for $p = 0$ (if estimated on the
maximum sample) matches the value reported in Table~1 of Lothian and Taylor
(1996) as $\tau_\mu$.

\item (DF test on the floating period) \enspace Apply the DF $t^\mu$-test to the floating exchange
rate period of 1974--1990. Because of the lagged dependent variable $z_{t-1}$, the
first equation is for $t = 1975$:
\[
z_{1975} = \text{const.} + \rho z_{1974} + \text{error}.
\]
So the sample size is~16. Can you reject the random walk null at 5 percent?

\item (DF test on the gold standard period) \enspace Apply the DF $t^\mu$-test to the gold standard
period of 1870--1913. Take $t = 1$ for 1871, so the sample size is~43. Can you
reject the random walk null at 5 percent? Is the estimate of $\rho$ (the $z_{t-1}$ coefficient in the regression of $z_t$ on a constant and $z_{t-1}$) closer to zero than in~(b)?

\item (Chow test) \enspace Having rejected the I(1) null, we proceed under the maintained
assumption that the log real exchange rate $z_t$ follows a stationary AR(1) process. To test the null hypothesis that there is no structural change in the AR(1)
equation, split the sample into two subsamples, 1791--1973 and 1974--1990,
and apply the Chow test. You may recall that the Chow test was originally
developed for models with strictly exogenous regressors. However, as you verified in Analytical Exercise~12 of Chapter~2, if the variables (the regressors and the dependent variable)
are ergodic stationary (which is the case here because $\{z_t\}$ is ergodic stationary when $|\rho| < 1$), then in large samples, with the sizes of both subsamples
increasing, $K \cdot F$ is asymptotically $\chi^2(K)$. This result does not require the
regressors to be strictly exogenous. Calculate $K \cdot F$ and verify the claim (see
pp.~498--502 of Lothian and Taylor, 1996) that ``there is no sign of a structural
shift in the parameters during the floating period.'' In the present case of a
lagged dependent variable, there is an issue of whether the equation for 1974,
\[
z_{1974} = \text{const.} + \rho z_{1973} + \text{error},
\]
should be included in the first subsample or in the second. It does not matter in
large samples; include the equation in the first subsample (so the sample size
of the first subsample is~183). (Note: $K \cdot F$ should be about 1.26.)

\item (The ADF-GLS) \enspace As was mentioned in Section~9.5, the ADF-GLS test achieves
substantially higher power at a cost of only slightly higher size distortions. The
ADF-GLS $t^\mu$- and $t^\tau$-statistics are calculated as follows. As in the ADF $t^\mu$- and
$t^\tau$-tests, the null process is
\[
y_t = d_t + z_t, \quad z_t = \rho z_{t-1} + u_t, \quad
\rho = 1, \quad \{u_t\} \text{ zero-mean I(0) process.}
\tag{3}
\]
Write $d_t = \mathbf{x}_t'\bm{\beta}$. For the case where $d_t$ is a constant, $\mathbf{x}_t = 1$ and $\bm{\beta} = \alpha$, while
for the case where $d_t$ is a linear time trend, $\mathbf{x}_t = (1, t)'$ and $\bm{\beta} = (\alpha, \delta)'$. We
assume the sample is of size $T$ with $(y_1, y_2, \ldots, y_T)$.

\begin{enumerate}[label=(\roman*)]
\item (GLS estimation of $\bm{\beta}$) \enspace In estimating the regression equation
\[
y_t = \mathbf{x}_t'\bm{\beta} + \text{error}, \quad (t = 1, 2, \ldots, T),
\tag{4}
\]
do GLS assuming that the error process is AR(1) with the autoregressive
coefficient of $\bar{\rho} = 1 + \bar{c}/T$. (The value of $\bar{c}$ will be specified in a moment.)
That is, estimate $\bm{\beta}$ by OLS by regressing
\[
y_1,\; y_2 - \bar{\rho} y_1,\; y_3 - \bar{\rho} y_2,\; \ldots,\; y_T - \bar{\rho} y_{T-1}
\]
on
\[
\mathbf{x}_1,\; \mathbf{x}_2 - \bar{\rho}\mathbf{x}_1,\; \mathbf{x}_3 - \bar{\rho}\mathbf{x}_2,\; \ldots,\; \mathbf{x}_T - \bar{\rho}\mathbf{x}_{T-1}.
\]
(Unlike in the genuine GLS, the first observations, $y_1$ and $\mathbf{x}_1$, are not
weighted by $\sqrt{1 - \bar{\rho}^2}$, but this should not matter in large samples.) Set
$\bar{c}$ according to
\[
\bar{c} = -7
\]
in the demeaned case (with $\mathbf{x}_t = 1$),
\[
\bar{c} = -13.5
\]
in the detrended case. (Roughly speaking, this choice of $\bar{c}$ ensures that the
asymptotic power at $c = \bar{c}$ where $\rho = 1 + c/T$ is nearly 50 percent which
is the asymptotic power achieved by an ideal test specifically tailored to
this alternative of $c = \bar{c}$.) Call the resulting estimator $\hat{\bm{\beta}}_{\text{GLS}}$ and define
$\tilde{y}_t \equiv y_t - \mathbf{x}_t'\hat{\bm{\beta}}_{\text{GLS}}$. (This should \emph{not} be confused with the GLS residual,
which is $(y_t - \bar{\rho}y_{t-1}) - (\mathbf{x}_t - \bar{\rho}\mathbf{x}_{t-1})'\hat{\bm{\beta}}_{\text{GLS}}$.)

\item (ADF without intercept on the residuals) \enspace Estimate the augmented autoregression \emph{without intercept} using $\tilde{y}_t$ instead of $y_t$:
\[
\tilde{y}_t = \rho\,\tilde{y}_{t-1} + \zeta_1 \cdot (\tilde{y}_{t-1} - \tilde{y}_{t-2})
+ \cdots + \zeta_p \cdot (\tilde{y}_{t-p} - \tilde{y}_{t-p-1}) + \text{error}.
\]
The $t$-value for $\rho = 1$ is the ADF-GLS $t$-statistic; call it the ADF-GLS $t^\mu$
if $\mathbf{x}_t = 1$ and the ADF-GLS $t^\tau$ if $\mathbf{x}_t = (1, t)'$. The limiting distribution
of the ADF-GLS $t^\mu$ is that of $DF_t$ (which is tabulated in Table~9.2, Panel~(a)), while that of the ADF-GLS $t^\tau$ is tabulated in Table~I.C.\ of Elliott et al.\
(1996). Reproducing their lower tail critical values
\[
-3.48 \text{ for } 1\%, \quad -3.15 \text{ for } 2.5\%, \quad
-2.89 \text{ for } 5\%, \quad -2.57 \text{ for } 10\%.
\]
Elliott et al.\ (1996) show for the ADF-GLS test that the asymptotic power
is higher than those of the ADF $\rho$- and $t$-tests, and that compared with
the ADF tests, the finite-sample power is substantially higher with only a
minor deterioration in size distortions.
\end{enumerate}

Apply the ADF-GLS $t^\mu$-test to the gold standard subsample 1870--1913.
Choose the lag length by BIC with $p_{\max}$ (the maximum number of lagged
changes in the augmented autoregression, denoted $p_{\max}(T)$ in the text) by~(9.4.31) (so $p_{\max} = 9$). Can you reject the null that the real exchange rate
is driftless I(1)? (Note: The BIC will pick 0 for $p$. The statistic should be
$-2.24$. You should be able to reject the I(1) null at 5 percent.)

\item (Optional analytical exercise on time aggregation) \enspace Suppose $y_t$ is a random
walk without drift. Create from this time-averaged series by the formula
\[
\bar{y}_1 = \frac{y_1 + y_2 + \cdots + y_n}{n}, \quad
\bar{y}_2 = \frac{y_{n+1} + y_{n+2} + \cdots + y_{2n}}{n}, \quad \ldots\;.
\]
(Think of $y_t$ as the exchange rate on day $t$ and $\bar{y}_j$ as an annual average of daily
exchange rates for year $j$.) Show that $\{\bar{y}_j\}$ is not a random walk. More specifically, show that
\[
\Corr(\bar{y}_{j+1} - \bar{y}_j,\; \bar{y}_{j+2} - \bar{y}_{j+1}) \;\dto\; 0.25 \quad \text{as } n \to \infty.
\]
\end{enumerate}

\end{problemset}

\begin{answers}

\subsection*{Analytical Exercises}

\answer{1}
Since $\E(\xi_0^2) < \infty$, $\xi_0/\sqrt{T}$ converges in mean square (and hence in probability)
to 0 as $T \to \infty$. Thus, the second term on the right-hand side of $(*)$ in the hint
can be ignored. Regarding the first term, we have
\[
\frac{\xi_T}{\sqrt{T}} = \frac{\xi_0}{\sqrt{T}}
+ \frac{\Delta\xi_1 + \Delta\xi_2 + \cdots + \Delta\xi_T}{\sqrt{T}}.
\]
The first term on the right-hand side of this equation can be ignored. By
Proposition~6.9, the second term converges to $N(0, \lambda^2)$ because $\Delta\xi_t$ is a linear
process with absolutely summable MA coefficients. So the first term on the
right-hand side of $(*)$ converges in distribution to $(\lambda^2/2)X$ where $X \sim \chi^2(1)$.
Since $\{\Delta\xi_t\}$ is ergodic stationary, the last term on the right-hand side of $(*)$
converges in probability to $\tfrac{1}{2}\,\E[(\Delta\xi_t)^2]$ by the Ergodic Theorem.

\answer{5}
The (1, 1) element of $\mathbf{A}_T$ converges to the random variable indicated in Proposition~9.2(c). Regarding the first element of $\mathbf{c}_T$, by Beveridge-Nelson,
\[
y_t = \psi(1)w_t + \eta_t + (y_0 - \eta_0), \quad w_t \equiv \varepsilon_1 + \varepsilon_2 + \cdots + \varepsilon_t.
\tag{1}
\]
Therefore,
\[
y_{t-1}^\mu = \psi(1) w_{t-1}^\mu + \eta_{t-1}^\mu,
\tag{2}
\]
where
\[
x_{t-1} = x_{t-1} - \bar{x}, \quad
\bar{x} = \frac{x_0 + \cdots + x_{T-1}}{T}
\quad \text{for } x = w, \eta.
\]
and
\[
\frac{1}{T}\sum_{t=1}^{T} y_{t-1}^\mu \cdot \varepsilon_t
= \psi(1)\frac{1}{T}\sum_{t=1}^{T} w_{t-1}^\mu \cdot \varepsilon_t
+ \frac{1}{T}\sum_{t=1}^{T}\eta_{t-1}^\mu \cdot \varepsilon_t.
\tag{3}
\]
From this we have
\[
\frac{1}{T}\sum_{t=1}^{T} y_{t-1}^\mu \cdot \varepsilon_t
\;\dto\; \tfrac{1}{2}\cdot\sigma^2 \cdot \psi(1) \cdot
\big\{[W^\mu(1)]^2 - [W^\mu(0)]^2 - 1\big\}.
\tag{4}
\]
The $\mathbf{A}_T$ matrix is asymptotically diagonal for the demeaned case too. The off-diagonal elements of $\mathbf{A}_T$ are equal to $1/\sqrt{T}$ times
\[
\frac{1}{T}\sum_{t=1}^{T}(\Delta y_{t-1})^{(\mu)}\, y_{t-1}^\mu.
\tag{5}
\]
Since $\sum_{t=1}^{T} y_{t-1}^\mu = 0$, this equals
\[
\frac{1}{T}\sum_{t=1}^{T}\Delta y_{t-1}\, y_{t-1}^\mu.
\tag{6}
\]
It is easy to show, by a little bit of algebra analogous to that in Review Question~3 of Section~9.4 and Proposition~9.2(d), that~(6) has a limiting distribution. So
$1/\sqrt{T}$ times~(6), which is the off-diagonal elements of $\mathbf{A}_T$, converges to zero
in probability. Then, using the same argument we used in the text for the case
without intercept, it is easy to show that $\hat{\zeta}_1$ is consistent for $\zeta_1$. So $\psi(1)$ is
consistently estimated by $1/(1 - \hat{\zeta}_1)$. Taken together,
\[
\frac{T \cdot (\hat{\rho}^\mu - 1)}{1 - \hat{\zeta}_1}
\;\dto\; DF_\rho^\mu.
\]

\subsection*{Empirical Exercises}

\answer{a}
With $p = 0$, the estimated AR(1) equation is
\[
z_t = \text{const.} + \underset{(0.0326)}{0.8869}\; z_{t-1}, \quad
SSR = 0.995135, \quad R^2 = 0.790, \quad T = 199.
\]
The value of $t^\mu$ is $-3.47$, which matches the value in Table~1 of Lothian and
Taylor (1996). The 5 percent critical value is $-2.86$ from Table~9.2, Panel~(b).

\answer{b}
The estimated AR(1) equation for 1974--1990 is
\[
z_t = \text{const.} + \underset{(0.1904)}{0.7704}\; z_{t-1}, \quad
SSR = 0.150964, \quad R^2 = 0.539, \quad T = 16.
\]
$t^\mu = -1.21 > -2.86$. We fail to reject the null at 5 percent.

\answer{c}
The estimated AR(1) equation for 1870--1913 is
\[
z_t = \text{const.} + \underset{(0.1045)}{0.7222}\; z_{t-1}, \quad
SSR = 0.069429, \quad R^2 = 0.538, \quad T = 43.
\]
$t^\mu = -2.66 > -2.86$. We fail to reject the null at 5 percent.

\answer{d}
The estimated AR(1) equation for 1791--1974 is
\[
z_t = \text{const.} + \underset{(0.0329)}{0.8993}\; z_{t-1}, \quad
SSR = 0.837772, \quad R^2 = 0.805, \quad T = 183.
\]
From this and the results in~(a) and~(b),
\[
K \cdot F = \frac{0.995135 - (0.837772 + 0.150964)}{\frac{0.837772 + 0.150964}{199 - 4}}
= 1.26.
\]
The 5 percent critical value for $\chi^2(2)$ is 5.99. So we accept the null hypothesis
of parameter stability.

\answer{e}
The ADF-GLS $t^\mu = -2.24$. The 5 percent critical value is $-1.95$ from Table~9.2, Panel~(a). So we \emph{can} reject the I(1) null.

\end{answers}
